% ================================================================
% SECTION 13: SOLVER SETUP
% ================================================================
\section{SolverSetup --- Orchestrator}
\label{sec:setup_solver}

\begin{center}
\code{bionetflux.setup\_solver} \quad(539 lines)
\end{center}

\code{SolverSetup} is the top-level orchestrator that 
ties together a problem module, an optional TOML configuration file, and
an optional \code{DomainGeometry} to produce all objects needed for a
simulation run.

% ------------------------------------------------------------------
\subsection{Constructor}

\begin{lstlisting}[language=Python, caption={SolverSetup constructor}]
class SolverSetup:
    def __init__(self,
                 problem_module: str = "bionetflux.problems.ooc_problem",
                 config_file: Optional[str] = None,
                 geometry: Optional[DomainGeometry] = None):
\end{lstlisting}

No work happens at construction time; the object stores the arguments
and defers all computation to \code{initialize()}.

% ------------------------------------------------------------------
\subsection{initialize()}

\begin{enumerate}
  \item Dynamically imports \code{problem\_module} via \code{importlib.import\_module}.
  \item Calls the module's \code{create\_global\_framework(geometry=\ldots, config\_file=\ldots)}.
  \item Stores the returned tuple:
    \begin{itemize}
      \item \code{problems}: list of \code{Problem} (one per domain).
      \item \code{global\_discretization}: \code{GlobalDiscretization}.
      \item \code{constraint\_manager}: \code{ConstraintManager}.
      \item \code{problem\_name}: human-readable string.
    \end{itemize}
  \item Sets \code{self.constraints = self.constraint\_manager} for backward compatibility.
\end{enumerate}

% ------------------------------------------------------------------
\subsection{Lazy Cached Properties}

The four heavy-weight objects are created only on first access and then cached:

\begin{longtable}{p{5cm}p{9.5cm}}
\toprule
\textbf{Property} & \textbf{What it creates} \\
\midrule
\code{elementary\_matrices} & \code{ElementaryMatrices(orthonormal\_basis=False)} \\
\code{static\_condensations} & One \code{StaticCondensationFactory.create()} per domain; matrices are built immediately. \\
\code{global\_assembler} & \code{GlobalAssembler.from\_framework\_objects(problems, global\_disc, static\_condensations, constraint\_manager)} \\
\code{bulk\_data\_manager} & \code{BulkDataManager(domain\_data\_list)} via \code{extract\_domain\_data\_list} \\
\bottomrule
\end{longtable}

% ------------------------------------------------------------------
\subsection{Solution-Vector Helpers}

\begin{longtable}{p{8cm}p{6.5cm}}
\toprule
\textbf{Method} & \textbf{Description} \\
\midrule
\code{create\_initial\_conditions() -> (List[np.ndarray], np.ndarray)} & Evaluate \code{problem.u0[eq](node)} for each domain; multipliers are zeros \\
\code{create\_global\_solution\_vector(trace\_solutions, multipliers) -> np.ndarray} & Concatenate domain traces + multipliers into one vector \\
\code{extract\_domain\_solutions(global\_solution) -> (List[np.ndarray], np.ndarray)} & Inverse of the above \\
\code{get\_problem\_info() -> dict} & Summary dict: name, DOF counts, time params, per-domain info \\
\code{validate\_setup(verbose=False) -> bool} & Round-trip test, bulk data test, residual/Jacobian test \\
\code{compute\_geometry\_from\_problems(name=None) -> DomainGeometry} & Build a \code{DomainGeometry} from problem extrema (fallback for legacy usage) \\
\bottomrule
\end{longtable}

% ------------------------------------------------------------------
\subsection{Module-Level Functions}

\begin{lstlisting}[language=Python, caption={Convenience factories}]
def create_solver_setup(problem_module, config_file=None,
                        geometry=None) -> SolverSetup:
    """Create and initialize a SolverSetup instance."""

def quick_setup(problem_module="bionetflux.problems.test_problem2",
                validate=True,
                config_file=None,
                geometry=None) -> SolverSetup:
    """Create, initialize, and optionally validate a SolverSetup."""
\end{lstlisting}

\code{quick\_setup} additionally calls \code{\_validate\_config\_compatibility}
when a \code{config\_file} is provided, checking that the TOML problem type
matches the problem module name.

% ------------------------------------------------------------------
\subsection{Backward-Compatible Shim}

A thin re-export shim at \code{src/setup\_solver.py} (top-level) allows legacy
\code{from setup\_solver import SolverSetup, quick\_setup} imports to work:

\begin{lstlisting}[language=Python, caption={src/setup\_solver.py}]
from bionetflux.setup_solver import (
    SolverSetup, create_solver_setup, quick_setup
)
\end{lstlisting}
